\begin{abstract} 
    Electrons were diffracted using a polycrystalline graphite lattice  
    with two different experimental setups. The wavelength of the electrons  
    was then calculated using de Broglie's postulate and their kinetic energy.  
    Finally, Bragg's law was applied to perform a linear fit and obtain  
    the interplanar distances of graphite.  

\end{abstract}
\begin{IEEEkeywords}
    Electron diffraction, de Broglie postulate, Bragg's law and crystalline structure of graphite.  
\end{IEEEkeywords}

\selectlanguage{spanish}
\begin{abstract}
    Se difractaron electrones mediante una red policristalina de grafito utilizando 
    dos montajes experimentales distintos. A partir de esto, se calculó la longitud 
    de onda de los electrones utilizando el postulado de de Broglie y su energía cinética. 
    Finalmente, se aplicó la ley de Bragg para realizar un ajuste lineal y obtener las 
    distancias interplanares del grafito. 

\end{abstract}
\begin{IEEEkeywords}
    Difracción de electrones, postulado de de Broglie, ley de Bragg y estructura cristalina del grafito.
\end{IEEEkeywords}
